% ID: SEL-D0076-OB-E03
% ESTADO: cotejada. Segunda pasada visual y borrador compilado revisados.
% Fuente: Selectividad/sel_2008_matematicas_aplicadas/Mat CCSS 4 - Examen y criterios de correción.pdf
% SHA-256: f1040adb7b9fb01c26219800e6212b9a26b93bb9d5e933743d18ee6ac26e163b
% Página: 2; opción B; ejercicio original 3.

\noindent\underline{ParteI}

Una caja contiene 12 bombillas, de las cuales 4 están fundidas. Se eligen, al azar y sin reemplazamiento, tres bombillas de esa caja.
\begin{enumerate}[label=\alph*)]
\item \textbf{(1 punto)} Calcule la probabilidad de que ninguna de las tres bombillas esté fundida.
\item \textbf{(1 punto)} Calcule la probabilidad de que las tres bombillas estén fundidas.
\end{enumerate}
\par\noindent\begin{minipage}{\linewidth}
\noindent\underline{Parte II}

El tiempo de utilización diaria de ordenador entre los empleados de una empresa sigue una distribución Normal de media $\mu$ y desviación típica $1.2$ horas.
\begin{enumerate}[label=\alph*)]
\item \textbf{(1.25 puntos)} Una muestra aleatoria de 40 empleados tiene una media del tiempo de utilización de $2.85$ horas diarias. Determine un intervalo de confianza, al $96\%$, para la media del tiempo de utilización diaria de ordenador.
\item \textbf{(0.75 puntos)} Calcule el tamaño mínimo que debería tener una muestra para estimar la media del tiempo de utilización diaria del ordenador con un error no superior a $0.75$ horas y el mismo nivel de confianza del apartado anterior.
\end{enumerate}
\end{minipage}\par
